\section{Structural reachability and reproducible state-space bounds}
\label{sec:state-space}

This section concerns boards after the mandatory spawn, including every
ordinary two-tile opening. Accumulated score and random-generator state
are not part of the counted object. We distinguish labeled boards from
orbits under the eight square symmetries. A run stopped at target $T$
counts boards whose maximum tile is strictly below $T$; individual losing
boards remain distinct and there are no artificial win or loss nodes.
The following arguments are human proofs with executable checks, not
extensions of the Lean formalization in Section~\ref{sec:lean}.

\subsection{A direct rank invariant}

For a $c$-cell board $B$, let $N_q(B)$ be the number of tiles with value
at least $2^q$.

\begin{theorem}[rank envelope]\label{ss:rank}
For every reachable board with $c\ge2$ cells,
\[
 N_q(B)\le c+2-q\qquad(2\le q\le c+2).
\]
Equivalently, decreasing nonzero ranks satisfy $r_p\le c+2-p$.
Consequently no tile exceeds $2^{c+1}$ and
\[
 \mass(B)\le 2^{c+2}-4.
\]
Equality in the mass bound requires the unique multiset
$\{2^{c+1},2^c,\ldots,4\}$.
\end{theorem}
\begin{proof}
The opening satisfies the inequalities; the case $q=2$ is the cell
capacity. Suppose the bounds hold before a slide and fix $q\ge3$.
Write $H=N_q(B)$ and let $j$ be the number of disjoint merges of two
rank-$(q-1)$ tiles. Only these merges can increase $N_q$.
If $j=0$, its old bound persists. Otherwise,
\[
 H+2j\le N_{q-1}(B)\le c+3-q,
 \qquad N_q(\text{after})\le H+j\le c+2-q.
\]
A spawn has rank at most 2, so the induction closes, including parallel
merges. No rank $c+2$ can occur, and larger ranks cannot arise without
first producing it. Sorting the threshold inequalities yields the rank
form. Sum its slotwise tile bounds to obtain the mass bound and its
unique equality multiset.
\end{proof}

This proof is independent of the earlier tight-moment argument. On
$4\times4$ it gives $\mass(B)\le262140$ and a tile bound of $131072$;
attainment still relies on the separately checked witnesses.

\subsection{Recent-spawn constraints}

Let $L_q(B)$ be the total mass in tiles strictly below $2^q$.
\begin{lemma}[recent-spawn mass]\label{ss:latency}
If $B$ contains a tile of rank at least $q\ge3$, then
\[
 L_q(B)\ge2(q-2).
\]
\end{lemma}
\begin{proof}
A spawn made $j$ rounds before $B$ has experienced at most $j$
subsequent merges along its ancestry, so its current descendant has
rank at most $j+2$. All of the latest $q-2$ spawns therefore have
descendants below rank $q$, and their disjoint ancestral masses sum to
at least $2(q-2)$. Reaching rank $q$ from an ordinary opening requires
at least $q-2$ rounds, so these spawns exist.
\end{proof}

A stronger necessary test records forced recent 4-spawns. Let $n_r$
count rank-$r$ tiles, let $K$ be the maximum rank, and let $E(B)$ be the
occupied cells whose deletion leaves occupancy compact toward some wall.
The last spawn must be an eligible 2 or 4. Set $f=0$ when an eligible 2
exists, otherwise $f=1$ when an eligible 4 exists, and reject if neither
does. For $r=3,\ldots,K$, set
\[
 j=r-2,\qquad U_r=\sum_{s<r}n_s2^{s-1}.
\]
Whenever $U_r\le j+f$, require $n_r>0$ and increase $f$ by one;
otherwise leave $f$ unchanged.

\begin{lemma}[causal supply]\label{ss:causal}
Every noninitial reachable board passes this test.
\end{lemma}
\begin{proof}
Let $F_j$ be the actual number of 4-spawns among the latest $j$ spawns.
Inductively $f\le F_j$. These $j$ spawns have descendants below rank
$r=j+2$ and occupy $2j+2F_j$ of the available lower-rank mass $2U_r$.
If $U_r\le j+f$, there is no remaining lower-rank mass for the next
older spawn to be a 2: its descendant would still be below rank $r$.
It must be a 4, and there must be a rank-$r$ tile to receive mass that
cannot fit below $r$. Increasing $f$ preserves the lower bound.
When the inequality fails, leaving $f$ unchanged is a relaxation.
A rejection can occur only at a missing rank $r\le K$, whence
$K\ge r+1=j+3$; at least $j+1$ rounds have then occurred, and the older
spawn used in the contradiction exists. No pre-opening spawn is assumed.
\end{proof}

For example, the inventory $\{128,8,2\}$ passes
Lemma~\ref{ss:latency} but fails Lemma~\ref{ss:causal}. The last spawn
must be the 2; the preceding spawn must be a 4 to contribute to the 8
with only one subsequent merge available. The latest five spawns would
then carry at least 12 mass below 128, but only 10 is present.

\subsection{Exact inverse slides}

\begin{theorem}[inverse line rule]\label{ss:inverse}
Let $(y_1,\ldots,y_m,0,\ldots,0)$ be a compact target line of length
$\ell$, with $y_i$ a power of two at least 2. Choose $s_i\in\{0,1\}$.
Use the input block $[y_i]$ if $s_i=0$, or $[y_i/2,y_i/2]$ if $s_i=1$,
the latter allowed only for $y_i\ge4$. These choices yield exactly all
preimages provided
\[
 m+\sum_i s_i\le\ell,\qquad
 s_i=0\ \Longrightarrow\ y_i\ne y_{i+1}/2^{s_{i+1}}
 \quad(i<m).
\]
Concatenate the blocks and insert zeros in any remaining positions.
\end{theorem}
\begin{proof}
Every output consumes either one input tile or two equal input tiles,
which gives the two block forms. A pair is consumed and cannot merge
onward in the same slide. A singleton must not equal the next block's
first tile, or the greedy scan would merge across the proposed boundary.
These conditions are necessary and make the scan consume exactly the
specified blocks, proving sufficiency. Inserting zeros preserves the
nonzero order.
\end{proof}

For a four-cell line the characterization is particularly simple.
A nonempty partial compact output is in the image exactly when its
nonzero sequence has no adjacent pair $2,2$; a full output is in the
image exactly when every adjacent pair differs. For unequal neighbors,
use a shifted sequence as a preimage. For a partial output $a,a,b$ split
the first $a$; for $a,b,b$ split the last $b$; for $a,a,a$ split the middle
$a$. The repeated value is at least 4 and Theorem~\ref{ss:inverse}
verifies every case. Full image lines and empty lines have only identity
preimages; every nonempty partial image line has a nonidentity preimage.

With $q=17$ allowed positive ranks, the image size within that rank range is
\[
 1+q+(q^2-1)+(q^3-2q+1)+q(q-1)^3=74{,}818.
\]
The checker exhausts all $18^4=104{,}976$ inputs and handles synthetic
rank-18 outputs separately. Board predecessors are obtained by deleting
each proposed last 2/4, choosing a common direction, and combining line
preimages, retaining nonidentity slides and post-spawn predecessors.
This describes immediate predecessors, not their reachability.

\subsection{A finite characterization and a numerical upper bound}

Let $I$ be the ordinary openings, $R$ the reachable boards, and $U$ a
proved finite superset of $R$. With exact successors define
\[
 F_0=U,\qquad F_{k+1}=I\cup\bigl(U\cap\operatorname{Succ}(F_k)\bigr).
\]
\begin{proposition}[rooted predecessor filtering]\label{ss:rooted}
The sequence decreases, contains $R$ at every stage, and stabilizes
at exactly $R$. For the $4\times4$ rank-bounded universe of mass at
least 4, $F_{131069}=R$.
\end{proposition}
\begin{proof}
Monotonicity and $F_1\subseteq U$ give descent. Reachable boards survive
because they are openings or have reachable predecessors. At a fixed
point every noninitial member has a predecessor in the same set.
Following predecessors decreases mass by 2 or 4, so it must end at an
opening. Every fixed-point member is therefore reachable. A chain in
the stated universe has at most $(262140-4)/2=131068$ edges; unrolling
the recurrence for 131069 stages leaves no orphan predecessor chain.
\end{proof}

We count a cheaper superset, using Theorem~\ref{ss:rank}, deletion-based
occupancy compaction, and Lemma~\ref{ss:causal}. For each symmetry,
occupied cells form cycles of its permutation, and invariant tile values
are constant on each cycle. Group cycles by length and eligibility,
choose low values 2 and 4, and assign ranks 3 through 17 with a finite
integer recurrence. At rank $r$, at most $18-r$ unassigned occupied
\emph{cells} may remain. Binomial factors choose the cycles; the running
lower-rank mass and $f$ implement the causal test. Saturation at 32 units
of value 2 is safe because every comparison threshold is below 32.

The fixed counts for the identity, quarter-turn, half-turn, three-quarter
turn, vertical reflection, horizontal reflection, and the two diagonals are
\[
\begin{gathered}
22{,}851{,}583{,}961{,}907{,}351{,}128;\quad9{,}616;\quad1{,}074{,}108{,}288;
\quad9{,}616;\\
1{,}296{,}452{,}906;\quad1{,}296{,}452{,}906;\quad
411{,}457{,}895{,}734;\quad411{,}457{,}895{,}734.
\end{gathered}
\]
Of the 480 labeled openings, 24 lie outside the compaction set. Their
fixed-count corrections are $(24,0,4,0,4,4,6,6)$. Burnside's lemma gives
\begin{align*}
 N_{\mathrm{labeled}}&\le22{,}851{,}583{,}961{,}907{,}351{,}152,\\
 N_{/D_4}&\le2{,}856{,}448{,}098{,}561{,}271{,}997.
\end{align*}
Separate Python-integer and JavaScript-BigInt implementations reproduce
all eight fixed counts and all opening corrections. These numbers are
exact sizes of necessary-condition supersets, not estimates or exact
counts of the reachable set. The exact inverse-slide value restrictions
are \emph{not} included in this numerical bound.

\subsection{Executed checks and reproducibility}

\begin{center}
\small
\begin{tabular}{@{}lrr@{}}
\toprule
Board and stopping rule & Symmetry classes & Labeled boards\\
\midrule
$2\times2$, continued &110&662\\
$2\times3$, continued &21{,}752&85{,}844\\
$3\times3$, continued &48{,}713{,}519&388{,}921{,}077\\
$4\times4$, stop at 8 &84{,}660&675{,}154\\
$4\times4$, stop at 16 &23{,}483{,}970&187{,}809{,}874\\
$4\times4$, mass at most 24 &67{,}104&535{,}148\\
\bottomrule
\end{tabular}
\end{center}

These are completed regression computations, not new record claims.
The separate raw-board Python search agrees with the entire symmetry
expansion of the C++ dumps for $2\times2$ and $2\times3$. On $2\times2$
the inverse relation is checked for all 1,296 syntactic target boards
through 32. Rooted filtering gives $896,686,666,662$ before stabilization.
The $3\times3$ benchmark is reproduced under the published convention;
mass-layer enumeration itself already appears in Lees-Miller's 2017
work~\cite{leesmiller2017b}, preceding the larger
$4\times3$ computation~\cite{kaneko2026}.

A 24-round $2\times2$ certificate ends on
\[
 \begin{pmatrix}0&0\\2&2\end{pmatrix}
 \longrightarrow
 \begin{pmatrix}2&8\\16&32\end{pmatrix},
 \qquad\operatorname{score}=180.
\]
It attains the ceiling without a final 4, demonstrating why both equality
cases in Theorem~\ref{thm:score} are required. Changing a final spawn
changes no merge and hence no score. The complete replay records every
move, spawn, intermediate board, and cumulative score.

Run \texttt{bash research/state-space/run.sh --large} from the repository
root. Each invocation uses a fresh directory, rebuilds the executable,
checks every reported aggregate against its layer table, compares both
bound implementations, and tests corrupt state dumps and invalid inputs.
The manifest records only executed cases, source hashes, tool versions,
exit codes, and output hashes. These hashes establish byte identity, not
completeness; the independent set and closure checks supply the latter
for the tested small boards. No full $4\times4$ cardinality, target-64/128
completion, or attainment of the $4\times4$ score ceiling is claimed.
